\subsection*{代价与限制}

为了支持动态形状，Croqtile 在编译时必须保守。这一限制让某些编译时检查做不了，只能挪到运行时。实现上已经尽量把这些检查放进主机代码来压低开销，但它们\textbf{仍然比固定形状的代码更保守}。程序员需要考虑这一限制，必要时添加用户级断言，加强代码安全。

\emph{匿名维度}和\emph{符号维度}都有一个重要限制：它们**只能用在参数列表里的 \emph{spanned data} 上**。内部缓冲区可以从输入派生形状，但不能用 \co@?@ 或新的符号名引入尚未指定的维度。例如，下面的代码会产生编译时错误：

\begin{lstlisting}[language=croqtile,caption={代价与限制}]
__co__ void foo(s32 [M] input) {
  s32 [M] derived;  // ok: the shape comes from the parameter
  s32 [?] fresh;    // error: an internal buffer cannot introduce an anonymous dimension
}
\end{lstlisting}
